% =====================================================================
\section{Cronograma de atividades executadas}
\label{sec:cronograma}

A Tabela~\ref{tab:cronograma} apresenta a distribuição das atividades ao longo da vigência.
As marcações correspondem aos períodos de execução; os bimestres anteriores a fevereiro de
2026 antecedem o registro versionado do projeto e devem ser conferidos contra o plano de
trabalho aprovado antes da entrega \pendente{conferir os dois primeiros bimestres com o plano
aprovado}.

\begin{table}[H]
\centering
\caption{Cronograma executado, por bimestre da vigência 2025--2026.}
\label{tab:cronograma}
\footnotesize
\setlength{\tabcolsep}{4pt}
\begin{tabular}{@{}p{6.6cm}cccccc@{}}
\toprule
\multirow{2}{*}{\textbf{Atividade}} & \multicolumn{6}{c}{\textbf{Bimestre}} \\
\cmidrule(l){2-7}
 & 1º & 2º & 3º & 4º & 5º & 6º \\
\midrule
Revisão bibliográfica em verificação formal, BMC/SMT e verificação de redes neurais
 & X & X & X &   &   &   \\
Preparação do ambiente e do repositório; estudo do ESBMC e de suas primitivas
 & X & X &   &   &   &   \\
Caso 1 --- redes quantizadas e \textit{stubs} de inferência
 &   & X & X & X &   &   \\
Caso 2 --- \textit{kernels} de inferência e estudo de escalabilidade
 &   &   & X & X &   &   \\
Caso 3 --- ciclo neuro-simbólico de geração e correção
 &   &   & X & X &   &   \\
Caso 4 --- controlador PID sob ruído e biblioteca de terceiros
 &   &   &   & X & X &   \\
Caso 5 --- blocos \textit{feed-forward} de modelos de linguagem
 &   &   &   & X & X &   \\
Caso 6 --- política de aprendizado por reforço
 &   &   &   & X & X &   \\
Caso 7 --- malha fechada \textit{cart-pole} com controlador DDPG
 &   &   &   &   & X & X \\
Construção da camada \texttt{core\_verify} e da suíte de regressão
 &   &   &   &   & X & X \\
Auditoria de evidências, correção das propriedades vácuas e reexecução dos
\textit{harnesses}
 &   &   &   &   &   & X \\
Pipeline IC3/PDR: geração do sistema de transição, síntese e verificação indutiva
 &   &   &   &   &   & X \\
Redação do artigo, do relatório final e preparação da apresentação
 &   &   &   &   & X & X \\
\bottomrule
\end{tabular}
\legend{\small Fonte: elaborada pelo autor.}
\end{table}

\subsection{Produtos gerados}

A Tabela~\ref{tab:produtos} relaciona os produtos entregues ao artefato correspondente no
repositório, de modo que cada item seja verificável de forma independente.

\begin{table}[H]
\centering
\caption{Produtos do ciclo e respectivos artefatos de verificação.}
\label{tab:produtos}
\footnotesize
\begin{tabular}{@{}p{6.0cm}p{8.0cm}@{}}
\toprule
\textbf{Produto} & \textbf{Artefato} \\
\midrule
Doze \textit{harnesses} de verificação reexecutáveis &
\texttt{scripts/gerar\_evidencias.py}, \texttt{results/logs/} \\
Tabela consolidada de evidências &
\texttt{results/EVIDENCIAS.md} \\
Invólucro programático com status tri-estado &
\texttt{core\_verify/esbmc\_caller.py} \\
Analisador de contraexemplos &
\texttt{core\_verify/SMT\_feedback\_parser.py} \\
Suíte de regressão, incluindo os quatro modos de falha &
\texttt{tests/}, com destaque para \texttt{tests/test\_ground\_truth.py} \\
Propriedades corrigidas da malha fechada e caracterização do envelope &
\texttt{cartpole/verify\_ddpg\_closed\_loop.py},
\texttt{cartpole/caracteriza\_prop\_b.py} \\
Pipeline de verificação ilimitada por IC3/PDR &
\texttt{ic3/gen\_transition\_system.py}, \texttt{ic3/validate\_forward.py},
\texttt{ic3/run\_pdr.sh} \\
Evidência das medições indutivas &
\texttt{ic3/EVIDENCIA.md} \\
Geradores e tabelas de consulta dos blocos FFN &
\texttt{cases/llm\_ffn\_verification/} \\
Regeneração reprodutível das figuras &
\texttt{artigo/figs/gerar\_figuras.sh} \\
Documentação de uso e de limitações conhecidas &
\texttt{pibic/README.md}, \texttt{cartpole/ESBMC\_NOTES.md} \\
Backlog auditado com procedência de evidência por tarefa &
\texttt{KANBAN.md} \\
\bottomrule
\end{tabular}
\legend{\small Fonte: elaborada pelo autor.}
\end{table}

\subsection{Participação em eventos e produção associada}

\pendente{listar aqui apresentação no Congresso de Iniciação Científica da UFAM, resumos
submetidos, participação em eventos e eventuais publicações decorrentes do plano}
